% ---------------------------------------------------------------------------
% 中文逐句主旨（与下方英文正文一一对应）
%
% Setup paragraph
% 1. Dense tiled convolution 本来就由 output-coordinate stream 驱动：output
%    coordinates 决定 input reads 与 output writes，而每个 position 内仍执行规则的
%    channel/kernel reduction。
% 2. WISEConv 利用这一性质，将 dense coordinate sequence 替换为 requested output
%    positions 的 worklist，从而把 position-level selectivity 接入现有的 tiled
%    convolution pipeline。
% 3. Worklist 可以从 schedule 中排除大量不需要的 output positions，并由此提高
%    useful-compute ratio，但减少 scheduled computation 本身并不保证 latency reduction。
% 4. 要将这种 selectivity 真正转化为 latency reduction，worklist 必须以足够低的成本
%    构造，并在 high useful-compute ratio 和 high throughput 下被消费。
% 5. WISEConv 因而围绕这些相互依赖的 performance requirements，共同设计
%    construction stage 与 compute stage，Figure overview 对此进行总结。
%
% Construction-stage paragraphs
% 1. 为避免 worklist construction cost 抵消 selective compute 节省的时间，WISEConv
%    让 construction 完全运行在 mask and coordinate metadata space，只读取 mask 并
%    生成 mask/coordinates，不访问 feature 或 weight tensors。
% 2. WISEConv 在一次 mask-space pass 中融合 Background Eq.~\ref{eq:mask} 的
%    output-mask propagation 与 tile-local compaction，同时生成 M_out 和 worklist W；
%    每个 spatial tile 的 requested coordinates 在 W 中形成一个按局部 dense traversal
%    order 排列的 segment，从而不经 global ordering 而保留 spatial grouping。
% 3. WISEConv 将 W 写入按静态 output-grid upper bound
%    |G|=H_out W_out 个 coordinates 预分配的 buffer，并由 device-resident counter 记录
%    worklist length n_W=|W|；construction 因而不依赖 host-visible active count 或 runtime
%    allocation。
% 4. 即使这样的 metadata-only pass 仍必须扫描当前 layer 的 mask grid，因此
%    compatible operators 继续传播并复用已有的合法 mask/worklist metadata，将一次
%    construction 的成本 amortize 到整个 compatible operator sequence。
% 5. Single-pass mask-space construction 降低每次 worklist construction 的成本，
%    cross-operator reuse 减少 construction 的发生次数，两者共同解决 Construction
%    Efficiency and Amortization。
%
% Compute-stage paragraphs
% 1. Worklist length 随 layer 和 input 变化，并且当前 n_W 只在 device 上可见；同步
%    获取这一 execution bound 会阻塞 per-frame path，但 workload decomposition 又必须
%    适应它，因为 padded work 和 exposed parallelism 分别影响 useful-compute ratio 与
%    GPU throughput。
% 2. 为在 host 不知道 execution bound 时完成当前工作，WISEConv 启动由 GPU
%    concurrent residency 和 dense output grid 的最大必要 work count 共同限定的
%    persistent thread blocks，并让这些 blocks 根据 device-resident n_W 以 grid-stride
%    方式消费 worklist，从而不需要 host synchronization。
% 3. 为在不查询精确 n_W 的情况下选择合适的 workload decomposition，WISEConv 在
%    deployment-time 对固定 envelope 的 DP、aligned Split-K 与 Hybrid Stream-K 做
%    length-indexed profiling，并把 policy table 放在 device 上；runtime 直接读取当前
%    device-resident worklist length，避免 input-activity prediction 和 host readback。
% 4. 因而，device-resident count 同时决定当前执行多少工作以及如何分解这些工作；
%    三种模式共享 reduction datapath，使 WISEConv 能够以 synchronization-free 的
%    per-frame path 消费 variable-length worklists，同时维持 high useful-compute ratio
%    和 high GPU throughput。
%
% Synthesis paragraph
% 1. Low-cost construction 与 efficient consumption 缺一不可：过高的 worklist
%    construction cost，或较低的 useful-compute ratio 与 compute throughput，都可能抵消
%    减少 scheduled computation 带来的收益。
% 2. 通过共同控制 worklist construction cost、useful-compute ratio 与 consumption
%    throughput，WISEConv 将 upstream spatial sparsity 转化为 end-to-end latency
%    reduction。
% ---------------------------------------------------------------------------

\section{\sysname Overview}
\label{sec:overview}

Dense tiled convolution already consumes a stream of output coordinates: each
coordinate determines the receptive-field reads and output write, while its
reduction over channels and kernel offsets remains regular. \sysname
exploits this interface by replacing the dense coordinate sequence with a
worklist of requested output positions, thereby introducing position-level
selectivity into the tiled pipeline. Removing unrequested positions from the
schedule raises the useful-compute ratio, but issuing fewer operations alone
does not guarantee lower latency. To translate selectivity into latency
reduction, the worklist must be inexpensive to construct and must be consumed with
both a high useful-compute ratio and high throughput. \sysname therefore
co-designs a construction stage and a compute stage around these interdependent
requirements, as Figure~\ref{fig:overview} summarizes.

% TODO: Replace this placeholder with the final two-stage overview figure.
\begin{figure*}[t]
  \centering
  \fbox{\parbox{0.9\textwidth}{\centering\vspace{2.5cm}%
    \textbf{[TODO]} \sysname overview. Left: the mask-space construction stage
    propagates the output mask, emits tile-local worklist segments into a
    preallocated buffer, and records the worklist length on the device;
    compatible operators reuse this metadata. Right: a persistent compute grid
    reads the device-resident worklist length and selects among data-parallel,
    split-$K$, and hybrid Stream-$K$ decompositions via a policy table
    calibrated offline.
    \vspace{2.5cm}}}
  \caption{The \sysname two-stage pipeline. Construction emits tile-locally
  ordered worklist segments using only mask and coordinate metadata, then reuses
  the metadata across compatible operators. Compute launches a persistent grid
  capped by SM occupancy and adapts its $K$-splitting decomposition on-device
  from the exact worklist length without host synchronization.}
  \Description{Placeholder for the \sysname overview figure. The construction
  stage emits and reuses mask and worklist metadata. The compute stage consumes
  a device-bounded worklist with a persistent grid and selects its
  decomposition mode from a device-resident policy table indexed by the current
  tile count.}
  \label{fig:overview}
\end{figure*}

\paragraph{Construction stage.}
To prevent worklist construction from erasing the savings of selective compute,
\sysname confines construction to mask and coordinate metadata, never reading
feature or weight tensors. In one mask-space pass, it fuses the output-mask
propagation in Eq.~\ref{eq:mask} with tile-local compaction, producing $M_{out}$
and a worklist $\mathcal{W}$; the requested coordinates from each spatial tile
form a segment in local dense-traversal order, preserving spatial grouping
without global ordering.
\sysname stores $\mathcal{W}$ in a buffer preallocated for
$|\mathcal{G}|=H_{out}W_{out}$ coordinates, the static upper bound on
$|\mathcal{W}|$. A device-resident counter records
$n_{\mathcal{W}}=|\mathcal{W}|$, so construction needs neither runtime
allocation nor a host-visible active count.

Even this metadata-only pass must scan the layer's mask grid. \sysname therefore
propagates and reuses the resulting mask and worklist metadata across compatible
operators instead of reconstructing it, amortizing one construction across the
compatible operator sequence. The fused pass lowers the cost of each
construction, while cross-operator reuse reduces how often it occurs; together,
they address the Worklist Construction Cost challenge.

\paragraph{Compute stage.}
The compute stage replaces only the coordinate source of a dense tiled
implicit-GEMM pipeline with the worklist; tensor layouts, weight access, and the
per-position channel-and-kernel reduction remain unchanged from the dense path.
The worklist length varies across layers and inputs and remains device-resident.
A fixed decomposition cannot serve the full range: too little parallelism
leaves the GPU underutilized at low activity, while unnecessary parallelism
adds overhead at high activity, both degrading throughput. Yet the conventional
approach of selecting a workload-size-specific decomposition on the host is
impractical, because reading the length back would incur a synchronization on
the critical path.

Because the worklist changes only the number of output tiles
while each tile's reduction datapath stays intact, \sysname launches a
persistent grid whose block count is set by the GPU's concurrent residency, not
by the current worklist length. The same blocks serve any worklist: when tiles
do not fill every block, idle blocks absorb reduction work from
undersubscribed tiles instead of sitting empty. Which blocks share and how much
each reduces is determined on-device: the kernel reads the exact worklist
length, looks up an offline-calibrated table keyed by that length, and
decomposes the workload accordingly. This on-device adaptation sustains GPU
throughput as the worklist length changes from frame to frame without any
host-side decision on the per-frame path.

Neither stage suffices alone: excessive construction cost, a low useful-compute
ratio, or low consumption throughput can each erase the benefit of issuing less
convolution work. By controlling all three jointly, \sysname converts upstream
spatial sparsity into end-to-end latency reduction.
